\documentclass[11pt]{article}
\usepackage{../../../tex/lecturestyle}

\title{Wavelet Matrix (Advanced)}
\author{}
\date{}

\begin{document}
\thispagestyle{fancy}

\lecturetitleband{Data Structures}{Wavelet Matrix (Advanced)}{7}

\section{Overview}
The companion lecture on the basics of the wavelet matrix showed how to
answer, in $O(H)$ time on a static array $A_0, \ldots, A_{N-1} \in [0,
2^H)$, the range $k$-th smallest and range count-in-$[a,b)$ queries, by
treating a query range $[L, R)$ as a node of a binary trie and walking it
top to bottom with an $O(1)$ child-range formula at every step. That same
walk generalizes a long way past those two queries. This lecture works
through the extensions that come up most often in contest problems:
XOR-shifted order statistics, range aggregates other than plain counting
(sums, and non-invertible aggregates like max), binary search against a
running sum, and reducing 2D rectangle queries to the same machinery via
coordinate compression. Assumes the basics lecture, plus familiarity with
segment trees and Fenwick trees.

\section{XOR-Shifted Order Statistics}

\begin{definitionbox}
Given $[L, R)$, $k$, and a mask $x \in [0, 2^H)$, find the $k$-th smallest
value of $A_i \oplus x$ over $i \in [L, R)$ (0-indexed).
\end{definitionbox}

Setting $x = 0$ is the plain range $k$-th-smallest query. The bit of $A_i
\oplus x$ at height $h$ equals the bit of $A_i$ at height $h$ when $x$'s
bit there is $0$, and its \emph{complement} when $x$'s bit is $1$. So the
walk is unchanged except at levels where $x$ has a $1$ bit: there, whatever
would normally be the ``left'' (bit-$0$) child of the node is really the
\emph{right} branch of the XOR-shifted ordering, and vice versa --- so swap
the two children's roles before comparing $k$ against the (swapped) left
child's size. Every step is still $O(1)$, so the whole query remains
$O(H)$.

\section{Aggregate Values Over a Range}

\subsection{Sums (or any group-valued aggregate)}
\begin{definitionbox}
Given $[L, R)$ and $[a, b)$, and an array $X_0, \ldots, X_{N-1}$ (not
necessarily equal to $A$), find $\sum_{i} X_i$ over $i \in [L, R)$ with
$a \le A_i < b$.
\end{definitionbox}

Exactly as range count-in-$[a,b)$ reduces to two count-less-than-$x$ calls,
range sum-in-$[a,b)$ reduces to two sum-less-than-$x$ calls:
$\mathrm{sumLess}(b) - \mathrm{sumLess}(a)$. To compute
$\mathrm{sumLess}(L, R, x)$, augment the wavelet matrix with, at each
height $h$, a prefix sum of $X$ restricted to the elements that go to the
$0$-branch at that height (in that height's pre-partition order) --- one
extra $O(N)$-length array per level, built alongside the existing rank
table during the same $O(NH)$ construction pass. The query walk is then
identical in shape to count-less-than: whenever $x$'s bit is $1$, add the
$0$-branch's $X$-sum over the current range to the answer and descend into
the $1$-branch; whenever it is $0$, descend into the $0$-branch with no
addition. Both $X = A$ (summing the values themselves) and an unrelated
attached array $X$ work identically --- nothing about the argument used
$X = A$. More generally this needs only that $X$ takes values in an
\emph{abelian group} (subtraction, i.e.\ the two-call trick, is what
requires invertibility).

\begin{theorembox}
Range sum-in-$[a,b)$, or any group-valued range aggregate, is answerable in
$O(H)$ time after $O(NH)$ preprocessing, by the same two-call
difference trick used for counting.
\end{theorembox}

\subsection{Non-invertible aggregates: max}
\begin{definitionbox}
Given $[L, R)$ and $[a, b)$, and an array $X$, find $\max_i X_i$ over
$i \in [L, R)$ with $a \le A_i < b$ (or report that no such $i$ exists).
\end{definitionbox}

The subtraction trick needs an inverse, and $\max$ has none: knowing
$\max(S \cup T)$ and $\max(T)$ does not recover $\max(S)$. A different
reduction is needed.

Think of the query as a search over trie nodes by \emph{value} interval
rather than by index. A node's value-interval $[c, d)$ (the root has
$[0, 2^H)$) either lies entirely inside $[a, b)$ (fold its aggregate into
the answer and stop), is disjoint from $[a, b)$ (prune it), or straddles the
boundary (split $[c,d)$ at its midpoint and recurse into both children,
narrowing $[l, r)$ to each child's index range via the same subtree-range
formula as before). This is precisely the recursion pattern of a segment
tree range query, and for the same reason it visits only $O(H)$ nodes in
total.

Realizing this needs a way to get the max of $X$ over an \emph{arbitrary}
index range at a fixed height $h$ --- not just a prefix, as the sum
trick needed. Build one auxiliary range-max structure (a segment tree
is the natural default, since $X$ can change) per height, over $X$
reordered into that height's array order. A query then costs $O(H)$ trie
nodes, each needing an $O(\log N)$ segment-tree query, for
$O(H \log N)$ total, or $O((N + Q)H \log N)$ across construction and $Q$
queries. Since the array is static, a sparse table or a sqrt tree can
replace the segment tree at each level if there is no need for updates
later, at the cost of $\Theta(NH\log N)$ space for a plain sparse table ---
watch memory here, since it is easy to multiply an already-large $NH$
structure by another $\log N$.

\subsection{Point updates to X}
If problems additionally need to change $X_i$ at a single index (leaving
$A$, and hence the trie structure itself, untouched), nothing about the
scheme above changes: a segment tree per height supports point updates
exactly as it supports the range query, so the same $O(H \log N)$ per
operation applies. When the aggregate is an invertible one (sum, XOR, any
abelian group), a Fenwick tree per height is a lighter-weight
substitute. Swapping which auxiliary structure backs each level is where
a lot of the wavelet matrix's flexibility in practice comes from.

\begin{notebox}
Changing $A_i$ itself (rather than an attached $X_i$) is a much harder
problem --- it would mean restructuring which leaf of the trie an element
belongs to, which cascades through every level. That is out of scope here;
a balanced BST-based structure is one way to support it, at real
implementation cost.
\end{notebox}

\section{Binary Search on a Running Aggregate}

\begin{definitionbox}
Given $[L, R)$ and a bound $S$, find the largest $x \le 2^H$ such that
$\sum_i A_i \le S$ over $i \in [L, R)$ with $A_i < x$.
\end{definitionbox}

The analogous question for \emph{counts} instead of sums is answered
directly by the range $k$-th-smallest query (find the value at rank $S$,
or just under it). For a running \emph{sum}, build the same per-level
prefix-sum tables as in the sum-aggregate case, then binary search is
unnecessary: walk the trie greedily from the root, and at each node decide
whether descending fully into the $0$-branch would keep the accumulated
sum at or below $S$. If yes, take the whole $0$-branch's contribution,
add it to a running total, and continue into the $1$-branch (since some
larger values may still fit within budget); if no, discard the $1$-branch
entirely and continue only into the $0$-branch. This greedy walk visits
$O(H)$ nodes and yields the same answer a binary search over $x$ would,
in $O(H)$ time instead of $O(H \log(2^H))$.

A close variant asks for the largest $x$ with
$\sum_{i} \max(x - A_i, 0) \le S$ over $[L, R)$ --- monotone in $x$, so
still solvable by walking down while separately tracking, for the value
range seen so far, both a running count and a running sum; the same
$O(H)$ greedy descent applies.

\section{Rectangles on a Point Set}

\subsection{Static points, counting}
\begin{definitionbox}
Given $N$ points $(A_i, B_i)$ and query rectangles $[L, R) \times [D, U)$,
count points inside each rectangle.
\end{definitionbox}

Sort the points by $x$-coordinate and re-index them by sorted position, so
that ``$L \le A_i < R$'' becomes a condition on a contiguous index range
(found by binary search on the sorted $x$-coordinates). Coordinate-compress
the $y$-coordinates to $[0, N)$ (or however many distinct values there
are). The problem is now exactly range count-in-$[a,b)$ on a static array
--- the compressed $y$-values, in sorted-by-$x$ order --- answerable by an
ordinary wavelet matrix of height $\lceil \log_2 N \rceil$. Construction
and $Q$ queries together cost $O((N + Q)\log N)$.

% FIGURE: wma_rect
\begin{figure}[H]
\centering
\begin{tikzpicture}[x=0.52cm,y=0.52cm]
\fill[lhi] (1.5,1.5) rectangle (7.5,5.5);
\draw[lectureink!15] (0,-0.5) -- (0,9.5);
\draw[lectureink!15] (-0.5,0) -- (9.5,0);
\draw[lectureink!15] (1,-0.5) -- (1,9.5);
\draw[lectureink!15] (-0.5,1) -- (9.5,1);
\draw[lectureink!15] (2,-0.5) -- (2,9.5);
\draw[lectureink!15] (-0.5,2) -- (9.5,2);
\draw[lectureink!15] (3,-0.5) -- (3,9.5);
\draw[lectureink!15] (-0.5,3) -- (9.5,3);
\draw[lectureink!15] (4,-0.5) -- (4,9.5);
\draw[lectureink!15] (-0.5,4) -- (9.5,4);
\draw[lectureink!15] (5,-0.5) -- (5,9.5);
\draw[lectureink!15] (-0.5,5) -- (9.5,5);
\draw[lectureink!15] (6,-0.5) -- (6,9.5);
\draw[lectureink!15] (-0.5,6) -- (9.5,6);
\draw[lectureink!15] (7,-0.5) -- (7,9.5);
\draw[lectureink!15] (-0.5,7) -- (9.5,7);
\draw[lectureink!15] (8,-0.5) -- (8,9.5);
\draw[lectureink!15] (-0.5,8) -- (9.5,8);
\draw[lectureink!15] (9,-0.5) -- (9,9.5);
\draw[lectureink!15] (-0.5,9) -- (9.5,9);
\draw[->,lectureink!70] (-0.5,-0.5) -- (10.2,-0.5) node[right,llabel] {position $i$};
\draw[->,lectureink!70] (-0.5,-0.5) -- (-0.5,10.2) node[above,llabel] {value};
\node[circle,fill=white,draw=lectureink!70,inner sep=1.6pt] at (0,3) {};
\node[circle,fill=white,draw=lectureink!70,inner sep=1.6pt] at (1,7) {};
\node[circle,fill=white,draw=lectureink!70,inner sep=1.6pt] at (2,1) {};
\node[circle,fill=lectureaccent,draw=lectureaccent,inner sep=1.6pt] at (3,5) {};
\node[circle,fill=lectureaccent,draw=lectureaccent,inner sep=1.6pt] at (4,2) {};
\node[circle,fill=white,draw=lectureink!70,inner sep=1.6pt] at (5,8) {};
\node[circle,fill=lectureaccent,draw=lectureaccent,inner sep=1.6pt] at (6,4) {};
\node[circle,fill=white,draw=lectureink!70,inner sep=1.6pt] at (7,0) {};
\node[circle,fill=white,draw=lectureink!70,inner sep=1.6pt] at (8,6) {};
\node[circle,fill=white,draw=lectureink!70,inner sep=1.6pt] at (9,3) {};
\node[llabel] at (4.5,-1.6) {points $(i,\,a_i)$ and the query rectangle};
\node[llabel,anchor=east] at (12.1,8.2) {array $a$};
\node[lcell,minimum width=0.52cm,minimum height=0.52cm,inner sep=0pt,font=\footnotesize\sffamily,] at (13.0,8.2) {3};
\node[llabel] at (13.0,9.2) {0};
\node[lcell,minimum width=0.52cm,minimum height=0.52cm,inner sep=0pt,font=\footnotesize\sffamily,] at (14.0,8.2) {7};
\node[llabel] at (14.0,9.2) {1};
\node[lcell,minimum width=0.52cm,minimum height=0.52cm,inner sep=0pt,font=\footnotesize\sffamily,lgray] at (15.0,8.2) {1};
\node[llabel] at (15.0,9.2) {2};
\node[lcell,minimum width=0.52cm,minimum height=0.52cm,inner sep=0pt,font=\footnotesize\sffamily,lhi] at (16.0,8.2) {5};
\node[llabel] at (16.0,9.2) {3};
\node[lcell,minimum width=0.52cm,minimum height=0.52cm,inner sep=0pt,font=\footnotesize\sffamily,lhi] at (17.0,8.2) {2};
\node[llabel] at (17.0,9.2) {4};
\node[lcell,minimum width=0.52cm,minimum height=0.52cm,inner sep=0pt,font=\footnotesize\sffamily,lgray] at (18.0,8.2) {8};
\node[llabel] at (18.0,9.2) {5};
\node[lcell,minimum width=0.52cm,minimum height=0.52cm,inner sep=0pt,font=\footnotesize\sffamily,lhi] at (19.0,8.2) {4};
\node[llabel] at (19.0,9.2) {6};
\node[lcell,minimum width=0.52cm,minimum height=0.52cm,inner sep=0pt,font=\footnotesize\sffamily,lgray] at (20.0,8.2) {0};
\node[llabel] at (20.0,9.2) {7};
\node[lcell,minimum width=0.52cm,minimum height=0.52cm,inner sep=0pt,font=\footnotesize\sffamily,] at (21.0,8.2) {6};
\node[llabel] at (21.0,9.2) {8};
\node[lcell,minimum width=0.52cm,minimum height=0.52cm,inner sep=0pt,font=\footnotesize\sffamily,] at (22.0,8.2) {3};
\node[llabel] at (22.0,9.2) {9};
\draw[decorate,decoration={brace,mirror,amplitude=3pt}] (14.5,7.55) -- (20.5,7.55) node[midway,below=3pt,llabel] {indices $[2,8)$};
\node[llabel,align=left,anchor=north west] at (12.5,5.6) {count the entries of this index range\\ whose value lies in $[2,6)$:\\ \textbf{3}};
\end{tikzpicture}
\caption{Counting points in a rectangle. Sorting the points by $x$ turns the query $[2,8)\times[2,6)$ into ``how many entries of the array $a_2,\dots,a_7$ lie in $[2,6)$'', a range-count query on a wavelet matrix; here the answer is 3.}\label{fig:wma_rect}
\end{figure}

\subsection{Points with insertions and deletions, offline}
When points are added and removed over time and queries ask for a maximum
attached value $X$ over a rectangle, read every query in advance. Each
point is then alive over a known, fixed time interval; replace its
insertion/deletion pair with ``its $X$ value applies during this interval,
and is the aggregate's identity element outside it.'' After
coordinate-compressing both spatial axes (and the discrete set of
event times), this is the max-with-point-updates construction from
earlier applied over an $O(N + Q)$-size compressed universe, giving
$O((N+Q)\log^2(N+Q))$ overall.

\section{Reference Implementations}
The examples below implement and cross-check (against brute force, over
every possible query) three representative techniques: XOR-shifted $k$-th
smallest, generalized range sum via an attached array, and 2D rectangle
counting by coordinate compression. The max-with-segment-tree-per-level,
point-update, binary-search, and offline-dynamic-points extensions above
follow the same pattern but are left as an implementation exercise, as the
source material for this curriculum itself suggests.

\codenote
\subsection{Rust}
\corelisting{Rust}{example.rs}

\subsection{Python}
\corelisting{Python}{example.py}

\section{Worked Example}
\begin{examplebox}
With $A = (3,1,2,1,1,0,3,2,2)$ and mask $x = 3$, the range $[2, 7)$ has
values $2,1,1,0,3$; XORing each with $3$ gives $1,2,2,3,0$, whose sorted
order is $0,1,2,2,3$. So the $0$-th smallest XOR-shifted value is $0$ and
the $4$-th is $3$ --- both confirmed by the brute-force cross-check in the
reference implementations, alongside every other $(l, r, k,
\mathrm{mask})$ combination on this array.
\end{examplebox}

\section{Further Reading}
\begin{itemize}[leftmargin=1.5em]
  \item AtCoder Algorithm Lectures index (source of this curriculum):
        \url{https://info.atcoder.jp/entry/algorithm_lectures/index}
  \item F. Claude, G. Navarro, A. Ordóñez Pereira, ``The wavelet matrix: An
        efficient wavelet tree for large alphabets'', \emph{Information
        Systems} 47, 2015, pp.\ 15--32.
\end{itemize}

\end{document}
